\documentclass[11pt,oneside,reqno]{amsart}
\usepackage[utf8]{inputenc}

\usepackage{enumitem}


\usepackage{amsthm}
\usepackage{mathtools}
\usepackage{amsmath}
\usepackage{amssymb}
\usepackage[all]{xy}
\usepackage{tikz}
\usepackage{tikz-cd}
\usepackage{mathtools}
\usepackage{mathrsfs}

\usepackage{color}

\def\rcomment#1{{\color{red}[{\sf #1}]}}
\def\bcomment#1{{\color{blue}[{\sf #1}]}}
\def\gcomment#1{{\color{green}[{\sf #1}]}}

%% blackboard bold

\newcommand\Z{{\mathbb Z}}
\newcommand\Q{{\mathbb Q}}
\newcommand\Qbar{{\overline{\mathbb Q}}}
\newcommand\N{{\mathbb N}}
\newcommand\R{{\mathbb R}}
\newcommand\C{{\mathbb C}}

\DeclareMathOperator{\sgn}{\text{sgn}}
\DeclareMathOperator{\per}{\text{Per}}
\DeclareMathOperator{\preper}{\text{Preper}}
\DeclareMathOperator{\tr}{tr}


%% bold face macros

\newcommand\br{{\mathbf r}}
\newcommand\bt{{\mathbf t}}
\newcommand\bu{\mathbf{u}}
\newcommand\bv{\mathbf{v}}
\newcommand\bw{\mathbf{w}}
\newcommand\bx{\mathbf{x}}
\newcommand\by{\mathbf{y}}
\newcommand\bX{\mathbf{X}}
\newcommand\bY{\mathbf{Y}}
\newcommand\bL{\mathbf{L}}

%%%%%%%%%%%%%%%%% environments %%%%%%%%%%%%%%%%%%

\numberwithin{equation}{section}

\newtheorem{theorem}{Theorem}[section]
\newtheorem{lemma}[theorem]{Lemma}
\newtheorem{prop}[theorem]{Proposition}
\newtheorem{cor}[theorem]{Corollary}

\theoremstyle{definition}
\newtheorem{definition}[theorem]{Definition}
\newtheorem{example}[theorem]{Example}
\newtheorem{problem}{Problem} 
\newtheorem{step}{Step}

\theoremstyle{remark}
\newtheorem{remark}[theorem]{Remark}
\newtheorem{notation}[theorem]{Notation}
\newtheorem{question}[theorem]{Question}
\newtheorem{conjecture}[theorem]{Conjecture}

\title{Lecture Notes: Complex Arithmetic and Geometry}
\author{S. Peter Jeong}

\begin{document}
\begin{abstract}
    A brief introduction to complex numbers, their arithmetic, and geometric representation.
\end{abstract}
\maketitle

\section{Complex Numbers} 
\subsection{Complex Arithmetic}
We begin with the definition of complex numbers and their arithmetic operations.
\begin{definition}
    A \textit{complex number} is an ordered pair of real numbers $(a,b)$, which we write as $a + bi$, where $i$ is the imaginary unit satisfying $i^2 = -1$. The set of all complex numbers is denoted by $\C$.
\end{definition}
Thus $\C$ extends the real numbers $\R$ using the imaginary unit $i$, allowing solutions to all polynomial equations. For example, the quadratic polynomial $x^2+2x+5$ has no real roots, 
but it has two complex roots $-1 + 2i$ and $-1 - 2i$.

For a complex number $z = a + bi$, the number $a$ is called the \textit{real part} of $z$, denoted $\operatorname{Re}(z)$, 
and the number $b$ is called the \textit{imaginary part} of $z$, denoted $\operatorname{Im}(z)$. Addition and multiplication of complex numbers 
do not deviate from the usual rules of algebra, as 
\[(a + bi) + (c + di) = (a + c) + (b + d)i\] and 
\[(a + bi)(c + di) = (ac - bd) + (ad + bc)i.\]
The quotient of two complex numbers is a complex number as well. If we assume $(a+bi)/(c+di) = (x+yi)$ given that $c+di \neq 0$,
we must have 
\[a+bi = (x+yi)(c+di) = (cx - dy) + (dx + cy)i\]
From this we obtain the system of equations
\begin{align*}
    a &= cx - dy \\
    b &= dx + cy
\end{align*}
which indeed has the unique solution 
\begin{align*}
    x &= \frac{ac + bd}{c^2 + d^2} \\
    y &= \frac{bc - ad}{c^2 + d^2}.
\end{align*}
Similarly, the square root of a complex number can be computed by solving 
\[(x + yi)^2 = a + bi\]
for some $x,y \in \R$. This leads to the system of equations
\begin{align*}
    x^2 - y^2 &= a \\
    2xy &= b.
\end{align*}

From these we obtain 
\[(x^2 + y^2)^2 = (x^2 - y^2)^2 + (2xy)^2 = a^2 + b^2\]
and thus 
\[x^2 + y^2 = \sqrt{a^2 + b^2}.\]
Combined with $x^2 - y^2 = a$, we have
\begin{align*}
    x^2 &= \frac{\sqrt{a^2 + b^2} + a}{2} \\
    y^2 &= \frac{\sqrt{a^2 + b^2} - a}{2}
\end{align*}
and both values are indeed nonnegative. Finally we have the general solution of the form
\[\sqrt{a + bi} = \pm \left( \sqrt{\frac{\sqrt{a^2 + b^2} + a}{2}} + i\frac{b}{|b|}\sqrt{\frac{\sqrt{a^2 + b^2} - a}{2}} \right).\]
given that $b \neq 0$. If $b = 0$, then the square root of $a$ is $\sqrt{a}$ if $a \geq 0$ and $i\sqrt{-a}$ if $a < 0$.


\subsection{Conjugation} Notice from the previous section that the only condition imaginary unit $i$ must satisfy is $i^2 = -1$. Since $i^2 = (-i)^2$, 
$-i$ must also satisfy any property that $i$ satisfies. The transformation $z = a + bi \mapsto \overline{z} = a - bi$ is called the \textit{complex conjugation}, and $\overline{z}$ is called the \textit{complex conjugate} of $z$.
The complex conjugation is an involution, meaning that $\overline{\overline{z}} = z$ for all $z \in \C$. It also satisfies the following properties:
\begin{gather*}
    \overline{z + w} = \overline{z} + \overline{w} \\
    \overline{zw} = \overline{z}\,\overline{w}.
\end{gather*} 
Consequently, if $ax = b$, then $\overline{a}\overline{x} = \overline{b}$, and thus $\overline{a/b} = \overline{a}/\overline{b}$ for $b \neq 0$. In general, 
if $R(z_1, z_2, z_3, \dots)$ is some rational operation on complex numbers $z_1, z_2, z_3, \dots$, then $\overline{R(z_1, z_2, z_3, \dots)} = R(\overline{z_1}, \overline{z_2}, \overline{z_3}, \dots)$. 
\begin{example}
    Consider the equation $c_nx^n + c_{n-1}x^{n-1} + \cdots + c_1x + c_0 = 0$ where $c_0, c_1, \dots, c_n \in \C$. If $z$ is a solution to this equation, then $\overline{z}$ is a solution to the equation
    \[\overline{c_n}x^n + \overline{c_{n-1}}x^{n-1} + \cdots + \overline{c_1}x + \overline{c_0} = 0.\] 
    In particular, if the coefficients $c_0, c_1, \dots, c_n$ are all real numbers, then $\overline{c_k} = c_k$ for all $k$, and it follows that the conjugate $\overline{z}$ is also a solution to the original equation.
    We thus have the familiar statement that the imaginary roots of a polynomial with real coeffcients always come in conjugate pairs. 
\end{example}

\subsection{Complex Numbers as Vectors} At this point, you may have noticed that the complex numbers are essentially pairs of real numbers that can be represented as points in the usual Cartesian plane as $(a, b)$, where $a$ is 
the $x$-coordinate (real axis) and $b$ is the $y$-coordinate (imaginary axis). And these points on the plane can always be represented as \textit{vectors} from the origin $(0,0)$ to the point itself using its 
\textit{magnitude} and the \textit{angle} it makes with the positive $x$-axis (measured counterclockedwise). This magnitude-angle representation can be explicitly found using basic trigonometry as follows. 

For a complex number $z = a + bi$, the magnitude of $z$ is given by the distance from the origin to the point $(a, b)$, which is
\[|z| = \sqrt{a^2 + b^2}.\]
This quantity is also called the \textit{absolute value} or \textit{modulus} of $z$. The angle $\theta = \tan^{-1}(b/a)$ that the vector from the origin to $(a, b)$ makes with the positive $x$-axis is called the \textit{argument} of $z$, denoted $\arg(z)$.
It follows that we can write $z$ in the form
\[z = (r,\theta)\]
where $r = |z|$ and $\theta = \arg(z) = \tan^{-1}(b/a)$. This is known as the \textit{polar form} of $z$. Conversely, given a polar coordinate $(r, \theta)$ of a complex number, 
we can write 
\begin{align*}
    a &= r\cos\theta \\
    b &= r\sin\theta.
\end{align*}
Hence we can write 
\[z = a+bi = r(\cos\theta + i\sin\theta).\]
This is known as the \textit{trigonometric form} of $z$. 
\begin{remark}
    Another way to compute the modulus is to use the complex conjugate. Notice that $z\overline{z} = (a + bi)(a - bi) = a^2 + b^2$, which is always nonnegative, and thus $|z|^2 = z\overline{z}$.
\end{remark}

\begin{remark}
    Division of complex numbers is massively simplified using the complex conjugate and modulus. For $z = a + bi$ and $w = c + di$ with $w \neq 0$, we have
\begin{align*}
    \frac{z}{w} &= \frac{a + bi}{c + di} \\
    &= \frac{(a + bi)(c - di)}{(c + di)(c - di)} \\
    &= \frac{(ac + bd) +    (ad - bc)i}{c^2 + d^2} \\
    &= \frac{ac + bd}{c^2 + d^2} + \frac{ad - bc}{c^2 + d^2}i \\
    &= \frac{z\overline{w}}{|w|^2}.
\end{align*}
\end{remark}

\begin{remark}
    Earlier we wrote $\arg(z) = \tan^{-1}(b/a)$, but since $\tan^{-1}\theta$ takes values in $(-\pi/2, \pi/2)$, the final value of $\arg(z)$ must be adjusted according to the quadrant in which $z$ lies. For example, if $a < 0$ and $b \geq 0$, then $\arg(z) = \tan^{-1}(b/a) + \pi$. If $a < 0$ and $b < 0$, then $\arg(z) = \tan^{-1}(b/a) - \pi$. If $a > 0$ and $b < 0$, then $\arg(z) = \tan^{-1}(b/a)$. If $a > 0$ and $b \geq 0$, then $\arg(z) = \tan^{-1}(b/a)$ as well. Finally, if $a = 0$ and $b > 0$, then $\arg(z) = \pi/2$, and if $a = 0$ and $b < 0$, then $\arg(z) = -\pi/2$. 
\end{remark}

\section{Inequalities and Identities} Now that we are familiar with both the arithmetic and geometric representation of complex numbers, we can derive some useful inequalities and identities.
\begin{lemma}\label{lemma:absineq}
    For any complex number $z$, we have $|\operatorname{Re}(z)| \leq |z|$ and $|\operatorname{Im}(z)| \leq |z|$.
\end{lemma}
\begin{proof}
    Let $z = a + bi$. Then 
    \[a \leq \sqrt{a^2 + b^2} = |z|.\]
    Similarly,
    \[b \leq \sqrt{a^2 + b^2} = |z|.\]
    Therefore,
    \[|\operatorname{Re}(z)| = |a| \leq |z| \quad \text{and} \quad |\operatorname{Im}(z)| = |b| \leq |z|.\]
\end{proof}

\begin{cor}[Cauchy-Schwarz Inequality for Complex Numbers]\label{cor:cauchy}
    For any two complex numbers $z$ and $w$, we have
    \[|\operatorname{Re}(z\overline{w})| \leq |z||w|\]
\end{cor}
\begin{proof}
    By Lemma \ref{lemma:absineq}, we have
    \[|\operatorname{Re}(z\overline{w})| \leq |z\overline{w}| = |z||\overline{w}| = |z||w|.\]
\end{proof}

Notice that we can rewrite Corollary \ref{cor:cauchy} as 
\[(x_1y_1 + x_2y_2)^2 \leq (x_1^2 + x_2^2)(y_1^2 + y_2^2).\]
for any real numbers $x_1, x_2, y_1, y_2$. In fact, this inequality extends to higher dimensions. If $x = (x_1, x_2, \dots, x_n)$ and $y = (y_1, y_2, \dots, y_n)$ are two $n$-dimensional vectors, then we have
\[(x_1y_1 + x_2y_2 + \cdots + x_ny_n)^2 \leq (x_1^2 + x_2^2 + \cdots + x_n^2)(y_1^2 + y_2^2 + \cdots + y_n^2).\]
This is the well-known Cauchy-Schwarz inequality for vectors in $\R^n$.

\begin{theorem}[Triangle Inequality]
    For any two complex numbers $z$ and $w$, we have
    \[|z + w| \leq |z| + |w|.\]
\end{theorem}
\begin{proof}
    We have
    \begin{align*}
        |z + w|^2 &= (z + w)(\overline{z} + \overline{w}) \\
        &= z\overline{z} + z\overline{w} + w\overline{z} + w\overline{w} \\
        &= |z|^2 + 2\operatorname{Re}(z\overline{w}) + |w|^2 \\
        &\leq |z|^2 + 2|z||w| + |w|^2 \\
        &= (|z| + |w|)^2.
    \end{align*}
    Taking the square root of both sides, we have $|z + w| \leq |z| + |w|$ as desired.
\end{proof}
The triangle inequality essentially states that on the complex plane, the shortest distance between two points is the straight line connecting them, and any other path connecting the two points must be at least as long as the straight line.
You can indeed verify this by simply drawing a few vectors. The triangle inequality also implies the following useful inequalities:
\begin{gather*}
    ||z| - |w|| \leq |z - w| \\
    |z| - |w| \leq |z + w| \leq |z| + |w|.
\end{gather*}

\begin{theorem}[Angle Addition of the Sine and Cosine]\label{thm:angleadd}
    For any two angles $\alpha$ and $\beta$, we have
    \begin{gather*}
        \sin(\alpha + \beta) = \sin\alpha\cos\beta + \cos\alpha\sin\beta \\
        \cos(\alpha + \beta) = \cos\alpha\cos\beta - \sin\alpha\sin\beta.
    \end{gather*}
\end{theorem}
\begin{proof}
    Left as an exercise.  
\end{proof}

An immediate result from the angle addition theorem above is that for two complex numbers $z_1 = r_1(\cos\theta_1 + i\sin\theta_1)$ and $z_2 = r_2(\cos\theta_2 + i\sin\theta_2)$, we have
\[z_1z_2 = r_1r_2(\cos(\theta_1 + \theta_2) + i\sin(\theta_1 + \theta_2)),\]
which implies 
\[\arg(z_1z_2) = \arg(z_1) + \arg(z_2).\]
Essentially, the product of two complex numbers is equivalent to some geometric transformation obtained by rotation and dilation. Futhermore, this kind of identity 
will be reoccuring in various forms throughout mathematics, such as the exponential and logarithmic rules 
\begin{gather*}
    b^{x+y} = b^xb^y \\
    \log(xy) = \log x + \log y.
\end{gather*}

Another important consequence of this geometric identity is as follows. 
\begin{theorem}[De Moivre's Theorem]
    For any angle $\theta$ and any integer $n$, we have
    \[(\cos\theta + i\sin\theta)^n = \cos(n\theta) + i\sin(n\theta).\]
\end{theorem}
\begin{proof}
    While this result is quite visual from the geometric perspective, we can prove this by induction on $n$. The base case $n = 0$ is trivial. For the inductive step, we have
    \begin{align*}
        (\cos\theta + i\sin\theta)^{n+1} &= (\cos\theta + i\sin\theta)^n(\cos\theta + i\sin\theta) \\
        &= (\cos(n\theta) + i\sin(n\theta))(\cos\theta + i\sin\theta) \\
        &= \cos((n+1)\theta) + i\sin((n+1)\theta).
    \end{align*}
\end{proof}
\end{document}